\chapter[\emj{💭}~DP - Memoization (Top-Down)]{\emj{💭}~DP - Memoization (Top-Down)}

\begin{dsaquees}
Un chef que anota recetas nuevas en un cuaderno 📖. 

Imagina que estás cocinando platillos complejos. La primera vez que
haces una salsa secreta, tardas mucho tiempo y esfuerzo. En lugar de
olvidar cómo la hiciste, anotas los pasos y el resultado en tu 
cuaderno.

La siguiente vez que alguien pide ese platillo, no vuelves a 
experimentar; simplemente abres tu cuaderno, lees la receta y 
¡listo! Tienes el resultado al instante. 

En programación, Memoization es el cuaderno donde guardas el resultado
de una función para ciertos argumentos, para no tener que volver a 
ejecutar toda la lógica si te vuelven a pedir lo mismo.
\end{dsaquees}
\begin{dsaotraforma}
Memoization es la técnica de DP basada en Recursión (Top-Down). 

Empiezas con el problema original (el más grande) y lo vas rompiendo 
en partes pequeñas. Antes de regresar el valor de una función, lo
guardas en una estructura de datos (generalmente un Array o un Hash Map)
donde la Key son los argumentos de la función.

Es muy natural de implementar si ya tienes una solución recursiva,
ya que solo requiere agregar el paso de "Revisar cuaderno" y 
"Anotar en cuaderno".
\end{dsaotraforma}
\begin{dsadiagrama}
\begin{verbatim}
Llamada: f(5)

1. ¿f(5) está en el cuaderno?
   - SÍ: Regresa el valor guardado. O(1) ⚡
   - NO:
     a. Calcula f(5) haciendo f(4) + f(3).
     b. Guarda el resultado: cuaderno[5] = resultado.
     c. Regresa el resultado.

¡Esto convierte un árbol de recursión gigante en un camino lineal!
\end{verbatim}
\end{dsadiagrama}
\begin{dsacomofunciona}
Para cualquier función recursiva:
1. Pasa una estructura `memo` como parámetro (o úsala global).
2. CASO BASE: Define cuándo detener la recursión.
3. REVISAR MEMO: Si el resultado ya está calculado, regrésalo.
4. CALCULAR Y GUARDAR: Calcula el valor, guárdalo en `memo` y regrésalo.
\end{dsacomofunciona}
\section*{PSEUDOCÓDIGO}
\hrulefill
\begin{verbatim}
función minMonedas(monto, monedas, memo):
    SI monto == 0: return 0
    SI monto < 0: return INFINITO
    SI monto EN memo: return memo[monto]

    minimo = INFINITO
    PARA cada moneda EN monedas:
        resultado = minMonedas(monto - moneda, monedas, memo)
        SI resultado != INFINITO:
            minimo = min(minimo, resultado + 1)

    memo[monto] = minimo
    return minimo
\end{verbatim}
\section*{CÓDIGO C++}
\hrulefill
\begin{lstlisting}
#include <iostream>
#include <vector>
#include <unordered_map>
#include <algorithm>
using namespace std;

// Problema: Coin Change (Mínimo de monedas para una cantidad)
// Tiempo: O(monto * n_monedas)
// Espacio: O(monto) para el memo
int solve(int amount, vector<int>& coins, unordered_map<int, int>& memo) {
    if (amount == 0) return 0;
    if (amount < 0) return -1;
    if (memo.count(amount)) return memo[amount];

    int min_coins = 1e9;
    bool possible = false;

    for (int coin : coins) {
        int res = solve(amount - coin, coins, memo);
        if (res >= 0) {
            min_coins = min(min_coins, res + 1);
            possible = true;
        }
    }

    return memo[amount] = (possible ? min_coins : -1);
}

int main() {
    vector<int> coins = {1, 2, 5};
    int amount = 11;
    unordered_map<int, int> memo;
    
    cout << "Mínimo de monedas para " << amount << ": " 
         << solve(amount, coins, memo) << endl; // 3 (5+5+1)
    
    return 0;
}
\end{lstlisting}
\begin{dsacomplejidad}
Aspecto        │ Con Memoization
  ───────────────┼────────────────────────────────────────────
  Tiempo         │ O(Estados Únicos * Trabajo por Estado)
  Espacio        │ O(Estados Únicos) + Stack de recursión

* Para Coin Change: O(Amount * N\_coins) tiempo. Mucho mejor que
  O(N\_coins \textasciicircum{} Amount) de la fuerza bruta.
\end{dsacomplejidad}
\begin{dsaentrevistas}

\end{dsaentrevistas}
\begin{dsaresumen}
• Top-Down: Del problema grande a los pequeños.
• Recursión + Diccionario/Array de resultados.
• Anota resultados para no repetir trabajo.
• Transforma complejidad exponencial en lineal/polinomial.
• Ideal para cuando la estructura del problema es un árbol.
════════════════════════════════════════════════════════════════
\end{dsaresumen}

